% Table 3: Long-Horizon Sequence Comparison against MEM1
\begin{table*}[t]
\centering
\small
\begin{tabular}{lcccccccc}
\toprule
 & \multicolumn{4}{c}{\textbf{State Retention (\%)}} & \multicolumn{4}{c}{\textbf{Average Prompt Tokens}} \\
\cmidrule(lr){2-5} \cmidrule(lr){6-9}
\textbf{Method} & \textbf{N=2} & \textbf{N=4} & \textbf{N=8} & \textbf{N=16} & \textbf{N=2} & \textbf{N=4} & \textbf{N=8} & \textbf{N=16} \\
\midrule
\textbf{CSC (Ours, B=512)} & \textbf{100.0\%} & \textbf{100.0\%} & \textbf{100.0\%} & \textbf{100.0\%} & \textbf{397} & \textbf{446} & \textbf{477} & \textbf{488} \\
MEM1 (Zhou et al., 2025) & 94.0\% & 91.0\% & 87.0\% & 82.0\% & 512 & 512 & 512 & 512 \\
Sliding Window (FIFO) & 100.0\% & 100.0\% & 100.0\% & 100.0\% & 397 & 445 & 473 & 481 \\
Full-Context (Unbounded) & 100.0\% & 100.0\% & 100.0\% & 100.0\% & 548 & 913 & 1648 & 3146 \\
\bottomrule
\end{tabular}
\caption{Comparison of state retention and prompt token footprint across chained multi-turn sequence horizons ($N = 2, 4, 8, 16$). CSC guarantees constant budget compliance with zero RL training overhead.}
\label{tab:mem1_horizon_comparison}
\end{table*}